\chapter{Evidências a incorporar na revisão final}\label{ap:pendencias}
Este apêndice reúne as pendências editoriais indicadas ao longo do texto. Ele funciona como apoio à preparação da versão final, sem substituir os resultados ou acrescentar procedimentos já executados.

\begingroup
\small
\setlength{\tabcolsep}{4pt}
\renewcommand{\arraystretch}{1.18}
\begin{longtable}{>{\raggedright\arraybackslash}p{\dimexpr0.23\textwidth-2\tabcolsep\relax}>{\raggedright\arraybackslash}p{\dimexpr0.35\textwidth-2\tabcolsep\relax}>{\raggedright\arraybackslash}p{\dimexpr0.42\textwidth-2\tabcolsep\relax}}
\caption{Evidências a incorporar na revisão final}\label{tab:apoio-1}\\
\toprule
Evidência & Destino no relatório & Informação necessária \\
\midrule
\endfirsthead
\multicolumn{3}{l}{\small Material de apoio à revisão (continuação)}\\
\toprule
Evidência & Destino no relatório & Informação necessária \\
\midrule
\endhead
\midrule
\multicolumn{3}{r}{\small Continua na próxima página.}\\
\endfoot
\bottomrule
\endlastfoot
Ambiente da versão final & Capítulo~\ref{cap:tcc2-4} & Data da avaliação, versão do motor, plataformas e configuração das máquinas. \\
Registro do teste integrado & Seção~\ref{sec:5.3.1} & Cenário, condições, resultado observado e vínculo com a build. \\
Rodadas de métricas & Seções 4.8 e 5.3.2 & Ambiente, participantes, duração, repetições e dados brutos. \\
Registro de CPU e memória, se utilizado & Seções 4.8 e 5.3.2 & Arquivo preservado e alinhamento temporal com as amostras do Godot. \\
Capturas de menu e resultado & Seções 5.1.3 e 5.2.9 & Incorporar as telas de entrada e encerramento; HUD, KOTH e buffs já ilustrados no texto. \\
Imagem da dashboard & Seção~\ref{sec:5.2.10} & Captura durante uma rodada identificada. \\
Build em checkout limpo & Seção~\ref{sec:5.2.11} & Preparação do ambiente, recursos necessários e execução dos artefatos. \\
Consolidação bibliográfica & Referências & Adaptação das entradas e citações ao formato exigido pelo curso. \\
\end{longtable}
\endgroup
